\subsection{Évolution des applications web}

Les applications de gestion d'entreprise ont connu une évolution majeure au cours des
deux dernières décennies. Les premières solutions reposaient sur des architectures
\textbf{monolithiques} : logique métier, interface et base de données regroupées dans
un même programme, souvent accessible via un navigateur avec rechargement complet de
page à chaque action.

\medskip

L'émergence des \textbf{API} (\textit{Application Programming Interface}) a séparé la
couche de présentation de la couche métier. Le navigateur ou une application mobile
consomme des services distants via \textbf{HTTP}, échangeant des données structurées en
\textbf{JSON}. Cette découpe a ouvert la voie aux architectures \textit{full-stack}
modernes, dans lesquelles un \textit{frontend} JavaScript riche dialogue avec un
\textit{backend} API indépendant~\cite{fielding2000rest}.

\subsection{Architecture client-serveur et tiers}

L'architecture \textbf{three-tier} (\textit{trois tiers}) structure l'application en
trois couches distinctes :

\begin{enumerate}
  \item \textbf{Couche présentation} : interface utilisateur (React), exécutée dans le
  navigateur.
  \item \textbf{Couche métier} : logique applicative, règles de gestion, authentification
  (Laravel).
  \item \textbf{Couche données} : persistance relationnelle (MySQL / PostgreSQL).
\end{enumerate}

Cette séparation améliore la maintenabilité : chaque couche peut évoluer indépendamment
tant que les contrats d'interface (API REST) sont respectés. Fowler~\cite{fowler2019patterns}
décrit ce découpage comme un patron fondamental des applications d'entreprise.

\begin{figure}[H]
  \centering
  \includegraphics[width=0.55\textwidth]{fig_2_1_three_tier_preview.png}
  \caption{Architecture three-tier de l'application CIVCO BTP}
  \label{fig:three_tier}
\end{figure}

\subsection{Single Page Application (SPA)}

Une \textbf{SPA} (\textit{Single Page Application}) charge une seule page HTML, puis
met à jour dynamiquement le contenu via JavaScript sans rechargement complet. Les
avantages pour une application métier sont significatifs :

\begin{itemize}
  \item \textbf{Fluidité} : navigation instantanée entre les modules (projets, clients,
  factures).
  \item \textbf{Expérience utilisateur} : interfaces réactives, proches des applications
  de bureau.
  \item \textbf{Réutilisation de l'état} : conservation du contexte utilisateur
  (filtres, session) entre les écrans.
\end{itemize}

Le couplage d'une SPA React avec une API Laravel constitue aujourd'hui un standard
industriel pour les tableaux de bord et outils de gestion internes.

\subsection{Panorama des stacks full-stack}

\begin{table}[H]
\caption{Panorama des stacks full-stack pour applications métier}
\label{tab:stack_panorama}
\centering
\small
\begin{tabular}{|l|l|p{8cm}|}
\hline
\textbf{Stack} & \textbf{Langages} & \textbf{Caractéristiques} \\
\hline
LAMP / LEMP & PHP + MySQL & Mature, hébergement abordable, Laravel/Symfony \\
\hline
MERN / MEAN & JavaScript & Un seul langage, Node.js + MongoDB \\
\hline
Django + React & Python + JS & Rapidité backend, ORM Django \\
\hline
\textbf{Laravel + React} & PHP + JS & API REST + SPA, écosystème riche \\
\hline
.NET + Angular & C\# + TS & Écosystème Microsoft, entreprise \\
\hline
\end{tabular}
\end{table}

Pour CIVCO BTP, la stack \textbf{Laravel + React} (Tableau~\ref{tab:stack_panorama})
offre le meilleur compromis entre productivité, courbe d'apprentissage et adéquation
avec les compétences visées en filière IIR : développement web, bases de données et
architecture logicielle.

\subsection{Échange de données et format JSON}

Les API modernes échangent des données structurées en \textbf{JSON} (\textit{JavaScript
Object Notation}), format léger et lisible, natif en JavaScript et bien supporté par
PHP/Laravel. Chaque ressource métier est sérialisée en objet JSON : un projet expose
son identifiant, son client associé, ses phases, son statut et ses dates ; un devis
expose ses lignes, montants et statut de validation.

\medskip

Ce format facilite le découplage frontend/backend : React consomme directement les
réponses JSON, les mappe vers ses composants et met à jour l'interface sans
transformation complexe. Les codes de statut HTTP (\texttt{200}, \texttt{201},
\texttt{422}, \texttt{403}) signalent le résultat de chaque opération de façon
standardisée.

\subsection{Patterns CRUD et ressources métier}

La majorité des modules de l'application suivent le pattern \textbf{CRUD}
(\textit{Create, Read, Update, Delete}) : créer un client, lister les projets,
modifier une facture, archiver un document. Laravel formalise ce pattern via les
contrôleurs de ressources et les routes RESTful, ce qui garantit une cohérence
d'API sur l'ensemble des modules (clients, projets, devis, factures, tâches,
contrats).
